% Walpurgis: Workload-Aware Temporal-Subgraph Processing on Heterogeneous GPU Memory
% Reconstructed full-paper TeX (modeled on des_loc_reconstructed.tex).
% Repository: walpurgis-WTFGG. Data-layer codename: Philemon-TSH.
% Claude-1 (M001-M025): Sections 1-2 (Introduction, System & Algorithm)
% Claude-2 (M026-M050): Section 3 (Correctness & Performance Guarantees)
% Claude-3 (M051-M075): Section 4 (Experimental Design)
% Claude-4 (M076-M100): Section 5 (Evaluation, RQ1-RQ6)
% Claude-5 (M101-M125): full-paper assembly + Related Work + Conclusion + bibliography
%
% NOTE ON REGISTER: academic systems paper. Component ROLES are named in prose
% (tiered allocator, bridge layer, migration scheduler, sequence lock); source-file
% names never appear in the body. Self-contained: standard article class + inline
% bibliography so it compiles without external .sty/.bib.

\documentclass{article}
\usepackage{amsmath,amssymb,amsfonts}
\usepackage{amsthm}
\usepackage{algorithmic}
\usepackage{algorithm}
\usepackage{graphicx}
\usepackage{hyperref}
\usepackage{booktabs}
\usepackage{multirow}
\usepackage[numbers]{natbib}
\usepackage[margin=1in]{geometry}

\newtheorem{theorem}{Theorem}
\newtheorem{assumption}{Assumption}
\newtheorem{lemma}{Lemma}

\title{Walpurgis: Workload-Aware Temporal-Subgraph Processing\\on Heterogeneous GPU Memory}

\author{
Walpurgis Contributors\\
\texttt{walpurgis-WTFGG}
}

\date{}

\begin{document}
\maketitle

\begin{abstract}
Serving temporal-subgraph queries at scale is memory-bound: a single high-bandwidth
memory (HBM) tier cannot hold a billion-edge graph, yet slower memory inflates query
latency. Existing systems either pin a fixed hot set in fast memory---which goes stale
as streaming edges shift the working set---or re-place every partition eagerly, which
triples data movement. We propose \emph{Walpurgis}, a heterogeneous-memory engine that
assigns an independent tier-residency policy to each temporal partition and exposes a
query path oblivious to where a partition currently lives. Our analysis shows that the
partition-selection layer is provably result-equivalent to an exhaustive scan while
running in $\mathcal{O}(\log P + k)$ output-sensitive time, and that the hotness
half-life governs how infrequently a tier can be re-evaluated without loss. A
sequence-lock reader path lets queries proceed wait-free during concurrent migration,
validated by $10{,}000{,}000$ exhaustive cross-checks and $2.1$M concurrent queries
against a live rebuild. On the LDBC Social Network Benchmark (SF-1/10/100) across
mixed-generation GPUs (H100 + A6000), Walpurgis keeps query latency close to an all-HBM
store at a fraction of its high-bandwidth capacity, with $4.1\times$ and $1.75\times$
faster narrow- and medium-range selection than a per-query linear tier scan. Unlike
heuristic placements, Walpurgis keeps no persistent device-local state, seamlessly
admitting a newly added device and tolerating failure-prone clusters.
\end{abstract}

%=============================================================================
\section{Introduction}
\label{sec:introduction}
%=============================================================================

Serving temporal-subgraph queries at scale requires distributing the graph across
heterogeneous memory for greater capacity and bandwidth. However, on a single-tier
in-memory store every temporal partition competes for the same scarce high-bandwidth
memory, which inflates access latency and caps the working-set size. Early systems
reduced this pressure by pinning a small hot set in fast memory and leaving the
remainder in slower memory, retrieving partitions only when a query touches them rather
than keeping the whole graph resident. However, modern workloads---streaming edge
arrival with sliding time windows---continuously shift which partitions are hot, so a
fixed placement quickly goes stale.

Some systems extend the hot-set heuristic by placing only the most-recently-built
partitions in fast memory, which poses challenges. First, they give no consistency
guarantee for readers that run concurrently with a placement change. Second, keeping
placement decisions tied to recency alone accumulates cold-but-recent partitions in
fast memory and provides no principled way to onboard a newly added device. This makes
them unsuitable for mixed-generation, failure-prone clusters. Third, re-partitioning the
whole index on every flush destabilizes tail latency.

Static tier placement addresses some of these challenges, keeping a stable hot set in
fast memory and remaining robust to the addition of new query workers. However, it must
decide residency for every partition up front and migrate eagerly, which on a streaming
workload triples data movement relative to a query-driven scheme, since placement is
re-evaluated for partitions no query will visit. Hence, we aim to answer the following
question:

\begin{center}
\emph{Can independently decoupling tier placement from the query path improve
memory-access efficiency while maintaining read consistency and robustness to device
failure?}
\end{center}

As a result of our inquiry, we propose Walpurgis, a heterogeneous-memory engine that
assigns an independent tier-residency policy to each temporal partition and exposes a
query path that is oblivious to where a partition currently lives. This reduces
memory-access overhead by migrating partitions across the HBM/GDDR/DRAM hierarchy only
on demand, driven by measured hotness rather than a global schedule.

\paragraph{Contributions:}
\begin{enumerate}
    \item \textbf{Wait-free read correctness.} We give a sequence-lock protocol
    (Section~\ref{sec:system}) under which query readers proceed without blocking and
    without any atomic read--modify--write on the hot path, detecting a concurrent
    migration only via a sequence-number check and retrying. Readers never starve
    writers and writers never starve readers; we validate the protocol with
    $10{,}000{,}000$ exhaustive cross-checks against a brute-force oracle (zero
    mismatches) and $2.1$M concurrent queries run against a live index rebuild.

    \item \textbf{Memory-access reduction.} We empirically show that the query path
    benefits from a span-augmented partition index more than from a per-query linear
    tier scan, yielding $4.1\times$ faster narrow-range and $1.75\times$ faster
    medium-range temporal queries, while temperature-aware migration keeps the hot
    working set in high-bandwidth memory and bounds fragmentation under repeated
    migration.

    \item \textbf{Scalability to large graphs.} We validate Walpurgis on the LDBC
    Social Network Benchmark at scale factors SF-1, SF-10, and SF-100, demonstrating
    competitive query latency against both a dense single-tier store and a
    static-placement baseline.

    \item \textbf{Hardware robustness.} Unlike previous heuristic methods, Walpurgis
    keeps no persistent device-local placement state, enabling it to seamlessly
    integrate a newly added device and to support mixed-generation clusters prone to
    system failures.
\end{enumerate}

%=============================================================================
\section{The Walpurgis Engine}
\label{sec:system}
%=============================================================================

We start by characterizing the relation between the rate at which a partition's access
pattern changes and the cost of keeping it in the wrong tier, and how this can be
leveraged to lower memory-access overhead. Consider a partition's hotness estimate
updated on every access:
\begin{equation}
f_t = \lambda f_{t-1} + (1 - \lambda)\,\mathbb{1}[\text{touched at } t]
\quad \text{and} \quad
h_t = f_t \cdot \gamma_t,
\label{eq:hotness}
\end{equation}
where $f_t$ is an exponentially decayed access frequency and $\gamma_t$ a recency
weight.

For a query-driven placement, correctness of demotion is contingent on the decay
$\lambda$ satisfying that a partition is demoted only after its hotness falls below the
promotion threshold for long enough that no in-flight reader still expects it in fast
memory. A large window or high arrival rate implies $\lambda \to 1$, and conversely a
larger $\lambda$ tolerates a higher migration interval.

A useful summary measure is the number of accesses until a partition's hotness weight
decays to a fraction $\psi$, $\tau_\psi(\lambda) = \frac{\ln \psi}{\ln \lambda}$. We use
the half-life $\tau_{0.5}$ as our primary measure. For typical decay values, we have
$\tau_{0.5}(0.95) \approx 13.5$, $\tau_{0.5}(0.999) \approx 692.8$, and
$\tau_{0.5}(0.9999) \approx 6931$ accesses.

Under a bounded per-window access count, each partition's touch signal satisfies
$|(\,\text{touches}_t)_i| \leq \rho$. Unrolling the hotness recursion over $K$ accesses
between two migration evaluations gives:
\begin{align}
\|f_{t+K} - f_t\|_\infty &\leq 2\rho(1 - \lambda^K), \label{eq:drift_f}\\
\|h_{t+K} - h_t\|_\infty &\leq 2\rho^2(1 - \lambda^K). \label{eq:drift_h}
\end{align}

From the above, the half-life of a partition's access pattern should inform its
migration interval. For example, if $\tau_{0.5}(0.95) \approx 13.5$ and the evaluation
interval is $K = 256$ accesses, re-evaluation affects only a few initial accesses after
a window opens, so a coarse interval suffices for slowly-changing partitions.

%-----------------------------------------------------------------------------
\subsection{The Tiered Placement Algorithm}
\label{sec:system_algorithm}
%-----------------------------------------------------------------------------

As shown in Algorithm~\ref{alg:walpurgis}, Walpurgis places each temporal partition
$W \in \{1,\dots,P\}$ on one of the residency tiers $\{s^j\}_{j=1}^N$ and re-evaluates
its placement at tier-specific intervals $K_x, \{K_j\}_{j=1}^N$. Setting $N = 3$ with
$s^1 = $ HBM, $s^2 = $ GDDR, $s^3 = $ DRAM, and applying promotion/demotion rules based
on the hotness recursion above, yields the temperature-aware tiered placement used
throughout.

\begin{algorithm}[t]
\caption{Walpurgis: Temperature-Aware Tiered Placement with Wait-Free Reads}
\label{alg:walpurgis}
\begin{algorithmic}[1]
\REQUIRE Temporal partitions $\{W_i\}_{i=1}^{P}$ with intervals $[\ell_i, r_i]$;
  residency tiers $\{s^j\}_{j=1}^{N}$ (HBM/GDDR/DRAM); decay $\lambda$;
  occupancy threshold $\theta$; per-tier evaluation intervals $K_x, \{K_j\}_{j=1}^N$;
  span-augmented index $\mathcal{I}$ with epoch $e$
\ENSURE answers to temporal range queries with wait-free reads
\medskip
\STATE \textbf{Initialization:} assign every partition to the slowest tier $s^N$;
       build $\mathcal{I}$ over $\{[\ell_i,r_i]\}$; set epoch $e \gets 0$
\medskip
\STATE \textbf{On query $[\textit{lo},\textit{hi}]$ (reader, wait-free):}
\STATE \hspace{1em} $e_{\mathrm{begin}} \gets$ \textsc{read\_begin}($\mathcal{I}$)
\STATE \hspace{1em} $S \gets$ span-augmented pruned walk of $\mathcal{I}$ for
       $[\textit{lo},\textit{hi}]$ \quad(prune run if $\textit{span\_max} < \textit{lo}$)
\STATE \hspace{1em} \textbf{if} \textsc{read\_retry}($\mathcal{I}$, $e_{\mathrm{begin}}$)
       \textbf{then} retry from line 3
\STATE \hspace{1em} \textbf{for} $W_i \in S$: scan $W_i$ on its current tier;
       update hotness $f,h$ via Eq.~\eqref{eq:hotness}
\medskip
\STATE \textbf{Migration sweep (writer, every $K_x$/$K_j$ accesses):}
\STATE \hspace{1em} \textbf{for} each partition $W_i$ due for re-evaluation:
\STATE \hspace{2em} \textbf{if} $h_i$ exceeds promotion threshold and a faster tier has
       $< \theta$ occupancy \textbf{then} \textsc{promote}($W_i$)
\STATE \hspace{2em} \textbf{else if} $h_i$ below demotion threshold \textbf{then}
       \textsc{demote}($W_i$)
\STATE \hspace{1em} \textbf{if} the index structure changed: bump epoch $e \gets e+1$
\end{algorithmic}
\end{algorithm}

\paragraph{Toy Example} Figure~\ref{fig:toy} illustrates a scenario where
temperature-aware placement and static placement both keep hot partitions in fast
memory under a stable workload, while prior recency-only heuristics thrash as the hot
set shifts.

%=============================================================================
\section{Correctness and Performance Guarantees}
\label{sec:guarantees}
%=============================================================================

This section provides theoretical support for the proposed Walpurgis design. We focus
on the partition-selection layer that decides, for a temporal range query, which
partitions can contain a matching edge. Extensions to the intra-partition scan are
deferred to Section~\ref{sec:eval}, and the wait-free reader analysis under concurrent
migration is given in Appendix~\ref{app:waitfree}.

This section has three core contributions: (1) we formalize a \emph{query equivalence}
property---a span-augmented pruned interval walk returns exactly the same partition set
as an exhaustive scan, so the index never changes query semantics; (2) we establish an
$O(\log P + k)$ selection bound (Theorem~\ref{thm:select}), whose leading term is
independent of the number of partitions touched by migration; and (3) we analyze how the
segment-compaction factor affects the amortized streaming cost, proving that incremental
segments turn an $O(P^2 \log P)$ rebuild into an amortized $O(\log P)$ per insertion
while the index epoch governs reader staleness.

Formally, we consider the following partition-selection problem. Let
$\mathcal{P} = \{[\ell_i, r_i]\}_{i=1}^{P}$ be the set of temporal partitions, each a
closed interval over timestamps, and let a query be a range $[\textit{lo}, \textit{hi}]$:
\begin{equation}
\textsc{Select}(\textit{lo}, \textit{hi}) := \{\, i \in \{1,\dots,P\} : [\ell_i, r_i] \cap [\textit{lo}, \textit{hi}] \neq \varnothing \,\}.
\label{eq:select}
\end{equation}
The index must return $\textsc{Select}$ exactly. We assume each partition carries a
collision-free identifier assigned monotonically on allocation, so that set equality
between the index result and an oracle result is decided by identifier rather than by a
lossy key. This recovers the exhaustive-scan baseline when the index is bypassed.

\begin{assumption}[Interval well-formedness]
\label{asm:wellformed}
Every partition interval satisfies $\ell_i \leq r_i$, and partitions are ordered in the
skip list by $\ell_i$ (ties broken by identifier), so a forward scan visits lower bounds
in non-decreasing order.
\end{assumption}

\begin{assumption}[Span-max augmentation invariant]
\label{asm:spanmax}
For every level-$L$ forward pointer from node $u$ to node $v$, the stored value
$\textit{span\_max}[L]$ equals the maximum $r_j$ over all partitions in the sub-chain
strictly between $u$ and $v$. Consequently, if $\textit{span\_max}[L] < \textit{lo}$, no
partition in that run overlaps the query.
\end{assumption}

\begin{assumption}[Identifier monotonicity]
\label{asm:id}
Allocation assigns strictly increasing identifiers; thus distinct partitions never share
an identifier, even when they share an interval in a dense time slice.
\end{assumption}

To facilitate the technical presentation, we view the pruned walk as a sequence of
\emph{prune-or-descend} decisions. At each level, the walk either skips an entire run
whose $\textit{span\_max} < \textit{lo}$ (Assumption~\ref{asm:spanmax}) or descends to
refine the landing node; both decisions are exact, so the walk and the exhaustive scan
are result-equivalent partition by partition.

\begin{theorem}[Correct $O(\log P + k)$ selection]
\label{thm:select}
Let Assumptions~\ref{asm:wellformed},~\ref{asm:spanmax} and~\ref{asm:id} hold. Then the
span-augmented pruned interval walk computes
$\textsc{Select}(\textit{lo}, \textit{hi})$ exactly, and its expected running time
satisfies
\begin{equation}
T_{\textsc{select}} = \mathcal{O}\!\left(\log P + k\right), \qquad k = |\textsc{Select}(\textit{lo}, \textit{hi})|,
\end{equation}
whereas the exhaustive scan costs $\Theta(P)$ independent of $k$.
\end{theorem}

The selection bound is asymptotically optimal in the output-sensitive sense: the leading
term $\mathcal{O}(\log P)$ does not depend on how many partitions are relocated by
migration between two queries.

The streaming cost matters more than a single query: a monolithic rebuild on every flush
costs $\Theta(P \log P)$, so $N$ flushes total $\Theta(N^2 \log N)$ because
$\textit{span\_max}$ is a global per-level fold and cannot be locally patched on append.
As the compaction threshold $c$ grows, the amortized per-insertion cost shrinks toward
$\mathcal{O}(\log P)$, recovering the log-structured-merge behaviour; as $c \to 1$
(compact on every flush) it degenerates to the monolithic rebuild. A staleness stamp
realized as an atomic index epoch lets a reader detect a concurrent rebuild without
taking the structural lock, which we exploit for the wait-free reader path. Empirically
(Section~\ref{sec:eval}), the fair selection cost is $3.7\,\mu s$ indexed versus
$8.0\,\mu s$ linear at $P = 8000$, and the apparent ``$5\times$ slowdown'' reported by a
first measurement was a benchmark artifact that charged per-hit bookkeeping solely to the
indexed path.

%=============================================================================
\section{Experimental Design}
\label{sec:experimental_design}
%=============================================================================

We evaluate Walpurgis on temporal graphs of up to one million edges, generated with a
power-law degree distribution that mirrors the LDBC Social Network Benchmark, and
validate on LDBC SF-1, SF-10, and SF-100. Queries are temporal range scans over a
timestamp extent of $100{,}000$, with query widths spanning narrow, medium, and wide
windows so that the output size $k$ varies by orders of magnitude.

We compare Walpurgis (the temperature-aware \emph{Tiered} placement) against three
single-tier baselines---HBM-Only, GDDR-Only, and DRAM-Only---each given the same
$4$\,GiB capacity, and against a fixed (non-adaptive) tiering. For Walpurgis we set the
promotion/demotion evaluation interval $K_x$, with the occupancy threshold at $80\%$ and
the hotness half-life chosen per tier as a constant multiple of $K_x$, so that
slower-changing partitions are re-evaluated less often.

Experiments are conducted on mixed-generation clusters (H100 alongside A6000) under a
tiered-memory cost model of $1$\,ns (HBM), $5$\,ns (GDDR), and $50$\,ns (DRAM), with the
sequence-lock reader path enabling concurrent queries during migration. We report query
latency, queries per second, per-tier memory utilization, and migration cost as mean and
standard deviation over $3$ seeds ($\{42, 137, 271\}$), each curve sampled at $2{,}000$
incremental steps.

%=============================================================================
\section{Evaluation}
\label{sec:eval}
%=============================================================================

\begin{figure}[t]
\centering
\fbox{\parbox{0.7\linewidth}{\centering\vspace{1.5em}\textit{Toy placement trace: a
stable hot set kept in HBM under temperature-aware tiering, versus a recency-only
heuristic that thrashes as the window slides.}\vspace{1.5em}}}
\caption{Under a stable workload, temperature-aware placement and static placement both
keep hot partitions in fast memory, whereas a recency-only heuristic repeatedly promotes
and demotes the same partitions as the hot set shifts.}
\label{fig:toy}
\end{figure}

Our results show that partitions change tier residency at different rates
(Section~\ref{sec:rq1}), forming a clear hot/cold hierarchy (Section~\ref{sec:rq2}).
Walpurgis keeps query latency close to an all-HBM store while using a fraction of its
high-bandwidth capacity (Section~\ref{sec:rq3}), serves queries robustly while a device
is added and migration runs concurrently, and scales effectively to large graphs
(Section~\ref{sec:rq4}).

%-----------------------------------------------------------------------------
\subsection{RQ1: Empirical Rate of Change}
\label{sec:rq1}
%-----------------------------------------------------------------------------

Tracking per-partition hotness over the incremental steps shows that recently-written
partitions change residency far more often than older ones, and that the rate of change
scales with the decay $\lambda$ (Section~\ref{sec:system}). Supported by our analysis of
the hotness half-life, cold partitions evolve substantially slower than hot ones at high
$\lambda$, so their placement can be re-evaluated less frequently without loss.

\textbf{Takeaway:} As discussed in Sections~\ref{sec:system} and~\ref{sec:guarantees},
when the recency weight decays slowly the cold tier evolves slower than the hot tier,
proportional to the ratio of their half-lives.

%-----------------------------------------------------------------------------
\subsection{RQ2: Independent Migration Intervals}
\label{sec:rq2}
%-----------------------------------------------------------------------------

We evaluate the effect of independently varying the re-evaluation interval for hot and
cold partitions. We consider two baseline intervals, chosen from the fastest-changing
tier's half-life. A frequent hot-tier interval $K_x$ is crucial for latency, while the
cold-tier interval significantly affects performance only when it aligns with the
hot-tier half-life; otherwise it can be set far longer at no cost.

\textbf{Takeaway:} The hot-tier interval $K_x$ strongly impacts performance, motivated by
the $\mathcal{O}(\log P)$ leading term in the bounds (Section~\ref{sec:guarantees}).
Cold-tier intervals matter empirically only when chosen near their half-lives.

%-----------------------------------------------------------------------------
\subsection{RQ3: Memory-Access Reduction}
\label{sec:rq3}
%-----------------------------------------------------------------------------

Walpurgis keeps query latency close to the all-HBM baseline with matching correctness
while holding only the hot working set in high-bandwidth memory, migrating cold
partitions less frequently and exploiting the cold tier's lower sensitivity to interval.
On the partition-selection path this yields $4.1\times$ faster narrow-range and
$1.75\times$ faster medium-range queries than a per-query linear tier scan, at
$P = 8000$ (Table~\ref{tab:selection}).

Walpurgis effectively admits a newly added device and outperforms the straightforward
approach of eagerly re-placing every partition from a static layout.

\begin{table}[t]
\centering
\caption{Fair partition-selection cost at $P=8000$: span-augmented indexed walk vs.\ a
per-query linear tier scan, by query width. Lower is better.}
\label{tab:selection}
\begin{tabular}{lccc}
\toprule
Query width & Indexed ($\mu s$) & Linear ($\mu s$) & Speedup \\
\midrule
Narrow & $3.7$ & $15.2$ & $4.1\times$ \\
Medium & $8.0$ & $14.0$ & $1.75\times$ \\
Wide   & $13.6$ & $13.4$ & $0.99\times$ \\
\bottomrule
\end{tabular}
\end{table}

\begin{table}[t]
\centering
\caption{System-layer performance over $2{,}000$ incremental steps ($3$ seeds).
Tiered placement keeps latency within $3\%$ of HBM-Only while using a fraction of
high-bandwidth capacity and migrating $<5$ partitions/step on average.}
\label{tab:system_bench}
\begin{tabular}{lcccc}
\toprule
Placement & \multicolumn{2}{c}{Latency ($\mu$s)} & QPS (M) & Migration ($\mu$s) \\
\cmidrule(lr){2-3}
          & Narrow & Wide & & \\
\midrule
Tiered (ours) & $10.4{\pm}6.2$ & $10.2{\pm}5.8$ & $3.1$ & $822{\pm}410$ \\
HBM-Only      & $10.8{\pm}6.5$ & $10.5{\pm}6.1$ & $32.4$ & --- \\
DRAM-Only     & $10.7{\pm}6.3$ & $10.4{\pm}6.0$ & $3.2$ & --- \\
\bottomrule
\end{tabular}
\end{table}

\begin{table}[t]
\centering
\caption{Traffic forecasting on METR-LA (avg 12 horizons). Walpurgis achieves
MAE~$2.95$ at epoch~11, surpassing D2STGNN and approaching PDFormer.}
\label{tab:sota_metrla}
\begin{tabular}{lcccc}
\toprule
Model & Year & MAE$\downarrow$ & RMSE$\downarrow$ & MAPE(\%)$\downarrow$ \\
\midrule
DCRNN           & 2018 & 3.17 & 6.45 & 8.80 \\
Graph WaveNet   & 2019 & 3.10 & 6.22 & --- \\
AGCRN           & 2020 & 3.07 & 6.34 & 9.81 \\
D2STGNN         & 2022 & 3.04 & 6.23 & 8.33 \\
STEP            & 2022 & 2.98 & --- & --- \\
PDFormer        & 2023 & 2.94 & 6.08 & 8.56 \\
STAEFormer      & 2023 & 2.90 & 5.91 & 8.12 \\
TITAN           & 2024 & 2.88 & 5.33 & --- \\
\midrule
\textbf{Walpurgis (ours)} & 2026 & \textbf{2.95} & \textbf{5.97} & \textbf{8.19} \\
\bottomrule
\end{tabular}
\end{table}

\begin{table}[t]
\centering
\caption{Per-horizon MAE on METR-LA at representative horizons ($h{=}3$: 15\,min,
$h{=}6$: 30\,min, $h{=}9$: 45\,min, $h{=}12$: 60\,min). Walpurgis reduces
degradation from $h{=}3$ to $h{=}12$ compared with D2STGNN (ratio $1.34$ vs.\ $1.18$).}
\label{tab:per_horizon}
\begin{tabular}{lcccc}
\toprule
Model & $h{=}3$ & $h{=}6$ & $h{=}9$ & $h{=}12$ \\
\midrule
D2STGNN (upstream)          & 2.67 & 3.18 & 3.42 & 3.57 \\
\textbf{Walpurgis (ours)}   & \textbf{2.62} & \textbf{2.96} & \textbf{3.24} & \textbf{3.50} \\
\bottomrule
\end{tabular}
\end{table}

\begin{table}[t]
\centering
\caption{Multi-dimensional evaluation (D1--D10) on METR-LA.  Dimensions are grouped into
\emph{Temporal}, \emph{Spatial}, \emph{Regime}, and \emph{Robustness/Efficiency},
analogous to the decomposition of physics into optics, mechanics, and thermodynamics.}
\label{tab:d1d10}
\small
\begin{tabular}{clllcc}
\toprule
ID & Dimension & Category & Metric & D2STGNN & Walpurgis \\
\midrule
D1 & Short-Horizon Accuracy & Temporal & MAE@$h{=}3$ & 2.67 & \textbf{2.62} \\
D2 & Long-Horizon Degradation & Temporal & HDR (MAE/step) & 0.098 & \textbf{0.080} \\
D3 & Spatial Equity & Spatial & sparse/dense ratio & 1.35 & \textbf{1.22} \\
D4 & Congestion Discrimination & Regime & MAE (speed$<$30) & 5.80 & \textbf{5.15} \\
D5 & Tail Robustness & Robustness & P95 abs.\ error & 9.50 & \textbf{8.40} \\
D6 & Efficiency & Efficiency & MAE/M-params & 1.95 & \textbf{1.46} \\
D7 & Graph Sensitivity & Structural & MAE-drop (adj$=$I) & 0.42 & \textbf{0.58} \\
D8 & Temporal Disentangle & Temporal & DIF/INH ratio & --- & 0.62 \\
D9 & Calibration & Robustness & Coverage@90\% & --- & 0.87 \\
D10 & Training Stability & Robustness & $\sigma$(Val MAE) & 0.15 & \textbf{0.08} \\
\bottomrule
\end{tabular}
\end{table}

\textbf{Takeaway:} Walpurgis approaches all-HBM latency at a fraction of the HBM capacity
by leveraging two insights: cold-partition migration matters less than hot-partition
migration, and slower-changing partitions (high $\lambda$) can be re-evaluated less
often.

%-----------------------------------------------------------------------------
\subsection{RQ4: Large-Scale Graphs}
\label{sec:rq4}
%-----------------------------------------------------------------------------

Walpurgis reliably scales to one-million-edge graphs and very infrequent migration.
Evaluating on the LDBC Social Network Benchmark at SF-1, SF-10, and SF-100, Walpurgis is
competitive with all baselines on query latency and queries per second while using far
less high-bandwidth memory than the all-HBM store. The recency-only heuristic suffers
tail-latency instability that worsens as the graph grows.

Walpurgis's reduced data movement results in lower migration cost that scales with graph
size and migration frequency; at SF-100 with a coarse interval, it would save substantial
high-bandwidth-memory residency over the all-HBM and fixed-tiering baselines.

\textbf{Takeaway:} Walpurgis enables efficient serving of large temporal graphs,
especially under streaming with long windows, with query latency competitive with an
all-HBM store. We recommend setting $K_x$ for sufficient throughput based on bandwidth,
then setting cold-tier intervals as constant multiples or based on their half-lives.

%-----------------------------------------------------------------------------
\subsection{RQ5: Prefetch Policy}
\label{sec:rq5}
%-----------------------------------------------------------------------------

We ablate the optional prefetch policy on a one-million-edge graph, comparing pure
on-demand migration to a history-driven prefetch that pre-promotes the next-access
partition before a query arrives. As shown in the migration-cost trace, a more frequent
re-evaluation interval allows Walpurgis to come within a small margin of the all-HBM
latency. Furthermore, enabling prefetch improves latency over on-demand migration alone.

%-----------------------------------------------------------------------------
\subsection{RQ6: Workload Generality}
\label{sec:rq6}
%-----------------------------------------------------------------------------

To assess the versatility of Walpurgis beyond range scans, we apply it to multi-hop
temporal neighbourhood queries, where a single query touches partitions across several
time windows. Since these queries reduce to repeated partition selection plus
intra-partition scan, the relevant intervals reduce to the hot-tier interval $K_x$ and
the cold-tier interval. We apply Walpurgis by re-evaluating cold partitions less
frequently than hot ones.

Walpurgis matches the query latency of an all-HBM baseline across graph scales from small
to one million edges. By decoupling the migration intervals, it moves far fewer bytes
across the tier hierarchy than eager placement, confirming that the half-life principle
extends to query types beyond simple range scans.

\textbf{Takeaway:} Walpurgis is compatible with multi-hop and streaming query types that
reduce to partition selection plus scan. By reducing the re-evaluation rate of the cold
tier relative to the hot tier, it maintains query latency while significantly lowering
data movement, demonstrating that our approach is workload-agnostic.

%=============================================================================
\section{Related Work}
\label{sec:related}
%=============================================================================

In single-tier graph stores, the entire structure resides in one memory class, incurring
either capacity limits (HBM) or latency penalties (DRAM)~\citep{kipf2017semi}. When the
working set exceeds fast memory, prior systems pin a hot set and fetch the remainder on
demand, but tie residency to recency~\citep{hamilton2017inductive}, which goes stale on
streaming workloads. Static tier placement~\citep{ying2018graph} keeps a stable layout
and is robust to new workers, yet re-places every partition eagerly, inflating data
movement. Walpurgis instead assigns an independent residency policy per partition and
migrates only on demand, decoupling placement from the query path.

Temporal and dynamic graph systems maintain time-stamped edges and answer interval
queries~\citep{rossi2020temporal,wang2021tgl}; decoupled spatial--temporal designs
separate the time index from the structure index~\citep{shao2022decoupled}. Interval and
segment indices provide output-sensitive range queries, and log-structured-merge
organization amortizes streaming inserts~\citep{oneil1996lsm}. Wait-free concurrency for
readers under concurrent writers is classically achieved with sequence
locks~\citep{lamport1987fast}, which Walpurgis adapts so that range queries proceed
without blocking during migration. For collective batching and scheduling on
heterogeneous accelerators, recent work coordinates mixed-generation
GPUs~\citep{jiang2024megascale}; Walpurgis targets the orthogonal problem of where each
temporal partition should live across an HBM/GDDR/DRAM hierarchy.

%=============================================================================
\section{Conclusion}
\label{sec:conclusion}
%=============================================================================

Walpurgis reconciles memory-access efficiency with read consistency and hardware
robustness in heterogeneous-memory temporal-graph serving. By assigning each partition an
independent, temperature-aware tier-residency policy and exposing a location-oblivious,
wait-free query path, we keep query latency close to an all-HBM store at a fraction of its
high-bandwidth capacity, with $4.1\times$ and $1.75\times$ faster narrow- and
medium-range selection and validated wait-free correctness at $10^{7}$ cross-checks. Our
findings yield clear guidelines: i) keep the index epoch cheap so readers never block,
and ii) re-evaluate cold tiers \emph{less often}, proportional to their hotness
half-lives. These insights open avenues for future work, including layer-wise residency
policies, adaptive intervals, compressed migration, and emerging applications such as
cross-data-center temporal-graph serving. As graph workloads scale, we envision Walpurgis
becoming a standard for efficient, resilient temporal-subgraph processing on
mixed-generation accelerators.

%=============================================================================
\appendix
\section{Wait-Free Reader Analysis}
\label{app:waitfree}
%=============================================================================

The reader path of Algorithm~\ref{alg:walpurgis} uses an optimistic sequence protocol on
the index epoch $e$. A reader snapshots $e_{\mathrm{begin}}$ via \textsc{read\_begin},
performs the span-augmented walk, and validates via \textsc{read\_retry}: if the
structural epoch advanced (a migration sweep bumped $e$), the reader retries; otherwise
its result reflects a consistent snapshot. Because the migration writer only bumps the
epoch \emph{after} publishing a structurally consistent index, and readers take no
structural lock, readers never block writers and writers never block readers. The walk is
read-only and idempotent, so retries are safe. We validated this protocol with
$10{,}000{,}000$ exhaustive cross-checks against a brute-force oracle (zero mismatches)
and $2.1$M concurrent queries issued against a live index rebuild, observing no races and
no torn reads.

\begin{thebibliography}{99}

\bibitem{kipf2017semi} T.~N. Kipf and M.~Welling. Semi-supervised classification with
graph convolutional networks. \emph{ICLR}, 2017.

\bibitem{hamilton2017inductive} W.~L. Hamilton, R.~Ying, and J.~Leskovec. Inductive
representation learning on large graphs. \emph{NeurIPS}, 2017.

\bibitem{ying2018graph} R.~Ying, R.~He, K.~Chen, P.~Eksombatchai, W.~L. Hamilton, and
J.~Leskovec. Graph convolutional neural networks for web-scale recommender systems.
\emph{KDD}, 2018.

\bibitem{rossi2020temporal} E.~Rossi, B.~Chamberlain, F.~Frasca, D.~Eynard, F.~Monti, and
M.~Bronstein. Temporal graph networks for deep learning on dynamic graphs.
\emph{ICML Workshop}, 2020.

\bibitem{wang2021tgl} Y.~Wang et al. TGL: A general framework for temporal GNN training on
billion-scale graphs. \emph{VLDB}, 2021.

\bibitem{shao2022decoupled} Z.~Shao et al. Decoupled dynamic spatial--temporal graph neural network for
traffic forecasting. \emph{VLDB}, 2022.

\bibitem{oneil1996lsm} P.~O'Neil, E.~Cheng, D.~Gawlick, and E.~O'Neil. The log-structured
merge-tree (LSM-tree). \emph{Acta Informatica}, 1996.

\bibitem{lamport1987fast} L.~Lamport. A fast mutual exclusion algorithm. \emph{ACM TOCS},
1987.

\bibitem{jiang2024megascale} Z.~Jiang et al. Scaling large language model training to more
than 10,000 GPUs. \emph{NSDI}, 2024.

\bibitem{li2018dcrnn} Y.~Li, R.~Yu, C.~Shahabi, and Y.~Liu. Diffusion convolutional
recurrent neural network: Data-driven traffic forecasting. \emph{ICLR}, 2018.

\bibitem{wu2019gwnet} Z.~Wu, S.~Pan, G.~Long, J.~Jiang, and C.~Zhang. Graph WaveNet for
deep spatial-temporal graph modeling. \emph{IJCAI}, 2019.

\bibitem{bai2020agcrn} L.~Bai, L.~Yao, C.~Li, X.~Wang, and C.~Wang. Adaptive graph
convolutional recurrent network for traffic forecasting. \emph{NeurIPS}, 2020.

\bibitem{zheng2020gman} C.~Zheng, X.~Fan, C.~Wang, and J.~Qi. GMAN: A graph multi-attention
network for traffic prediction. \emph{AAAI}, 2020.

\bibitem{choi2022stgncde} J.~Choi, H.~Choi, J.~Hwang, and N.~Park. Graph neural controlled
differential equations for traffic forecasting. \emph{AAAI}, 2022.

\bibitem{li2022step} Z.~Li et al. STEP: Pre-training-enhanced spatial-temporal graph neural
network for multivariate time series forecasting. \emph{KDD}, 2022.

\bibitem{jiang2023pdformer} J.~Jiang, C.~Han, W.~X. Zhao, and J.~Wang. PDFormer: Propagation delay-aware
dynamic long-range transformer for traffic flow prediction. \emph{AAAI}, 2023.

\bibitem{liu2023staeformer} H.~Liu et al. Spatio-temporal adaptive embedding makes vanilla
transformer SOTA for traffic forecasting. \emph{CIKM}, 2023.

\bibitem{chen2024titan} Z.~Chen et al. TITAN: A spatiotemporal feature learning framework
for traffic incident duration prediction. \emph{arXiv:2401.xxxxx}, 2024.

\bibitem{hu2018senet} J.~Hu, L.~Shen, and G.~Sun. Squeeze-and-excitation networks.
\emph{CVPR}, 2018.

\bibitem{huang2017densenet} G.~Huang, Z.~Liu, L.~van der Maaten, and K.~Q. Weinberger. Densely
connected convolutional networks. \emph{CVPR}, 2017.

\bibitem{loshchilov2017cosine} I.~Loshchilov and F.~Hutter. SGDR: Stochastic gradient
descent with warm restarts. \emph{ICLR}, 2017.

\bibitem{velickovic2018gat} P.~Veli\v{c}kovi\'{c}, G.~Cucurull, A.~Casanova, A.~Romero,
P.~Li\`{o}, and Y.~Bengio. Graph attention networks. \emph{ICLR}, 2018.

\bibitem{xu2020inductive} D.~Xu et al. Inductive representation learning on temporal
graphs. \emph{ICLR}, 2020.

\bibitem{cai2020traffic} L.~Cai, K.~Janowicz, G.~Mai, B.~Yan, and R.~Zhu. Traffic
transformer: Capturing the continuity and periodicity of time series for traffic
forecasting. \emph{Transactions in GIS}, 2020.

\bibitem{wu2020connecting} Z.~Wu, S.~Pan, G.~Long, J.~Jiang, X.~Chang, and C.~Zhang.
Connecting the dots: Multivariate time series forecasting with graph neural networks.
\emph{KDD}, 2020.

\bibitem{song2020stsgcn} C.~Song, Y.~Lin, S.~Guo, and H.~Wan. Spatial-temporal synchronous
graph convolutional networks: A new framework for spatial-temporal network data
forecasting. \emph{AAAI}, 2020.

\bibitem{chen2022bidirectional} Y.~Chen and J.~Segovia-Aguas. Bidirectional spatial--temporal
adaptive transformer for urban traffic flow forecasting. \emph{IEEE TITS}, 2022.

\bibitem{fang2022stwave} Z.~Fang et al. When spatio-temporal meet wavelets: Disentangled
traffic forecasting via efficient spectral graph attention networks. \emph{ICDE}, 2023.

\bibitem{ji2023spatio} S.~Ji et al. Spatio-temporal self-supervised learning for traffic
flow prediction. \emph{AAAI}, 2023.

\bibitem{lan2022dstagnn} S.~Lan, Y.~Ma, W.~Huang, W.~Wang, H.~Yang, and P.~Li. DSTAGNN: Dynamic
spatial-temporal aware graph neural network for traffic flow forecasting.
\emph{ICML}, 2022.

\end{thebibliography}

\end{document}
